\section{Lean formalization outline and path to object iPRA}\label{sec:formalization}

This section is deliberately more concrete than a list of desired theorems.  It
specifies representations, module boundaries, induction invariants and the
final theorem type.  The accompanying \code{lean/} directory mirrors this
plan.  Those files are a scaffold: the artifact-building environment did not
contain Lean, and the analytic object-theory derivations have not yet been
implemented.

\subsection{The endpoint that counts}

The target is not merely a theorem about Lean's \code{Nat}, \code{Rat}, or a
host-level complex type.  At the pinned ProveIt commit, the endpoint is
schematically
\begin{lstlisting}[language={}]
namespace LeanProofs.ConciseIPRA

-- phi : Fm PTm 0
example : [] |- [PRAAx] phi := by
  ...  -- finite object-language derivation

end LeanProofs.ConciseIPRA
\end{lstlisting}
The repository notation is \(\code{[]}\vdash_{\code{PRAAx}}\varphi\).  The syntax
\code{PR} has constructors for zero, successor, projection, composition and
primitive recursion.  \code{PRAAx} contains the defining equations, successor
axioms and quantifier-free induction, while \code{PRA.qfInd} is the convenient
induction constructor.  After an ordinary derivation is obtained, the
\code{pra\_iff\_native} bridge transfers it to the concise native scalar core.

A host theorem is useful during development but is not the final result.  Every
module should therefore have three layers:

\begin{enumerate}
\item executable host definitions and regression theorems;
\item graph formulas and bounded checkers in \code{PTm}/\code{Fm};
\item a theorem whose type is \code{Prf PRAAx [] phi}.
\end{enumerate}

\subsection{Module graph}

A practical dependency order is
\begin{verbatim}
BarnesG/FiniteProducts
  -> BarnesG/IntegerValues
  -> BarnesG/RationalBalls
  -> BarnesG/ElementaryNames
  -> BarnesG/CanonicalFactor
  -> BarnesG/CanonicalProduct
  -> BarnesG/FunctionalEquation
  -> BarnesG/Taylor
  -> BarnesG/Integral
  -> BarnesG/Derivative
  -> BarnesG/Reflection
  -> BarnesG/Multiplication
  -> HurwitzZeta/DerivativeMinusOne
  -> BarnesG/Asymptotic
  -> BarnesG/SpecialValues
  -> BarnesG/Characterization
  -> BarnesG/ObjectTargets
\end{verbatim}
The order is intentionally not the order of the source pages.  It minimizes
circularity and permits early object-level milestones.

\subsection{Stage 1: exact finite core}

Introduce primitive-recursive symbols:
\[
 \operatorname{fac},\quad g,\quad \operatorname{sf},\quad
 \operatorname{hyp},\quad \operatorname{pow},\quad
 \operatorname{facBlock},\quad c,\quad \operatorname{oddDF}.
\]
The intended equations are
\begin{align*}
 \operatorname{fac}(0)&=1,&
 \operatorname{fac}(n+1)&=(n+1)\operatorname{fac}(n),\\
 g(0)&=1,&g(n+1)&=\operatorname{fac}(n)g(n),\\
 \operatorname{hyp}(0)&=1,&
 \operatorname{hyp}(n+1)&=\operatorname{hyp}(n)(n+1)^{n+1},\\
 c(0)&=1,&c(k+1)&=c(k)\operatorname{oddDF}(k+1).
\end{align*}
All are direct \code{PR.prec} definitions.  Prove:
\begin{align}
 g(n+1)&=\operatorname{sf}(n),\label{eq:lean-sf-target}\\
 g(n)\operatorname{hyp}(n)&=\operatorname{fac}(n)^n,
 \label{eq:lean-hyp-target}\\
 g(k+n)&=g(k)\operatorname{facBlock}(k,n).
 \label{eq:lean-block-target}
\end{align}
The first and third inductions are telescopes.  The second uses the invariant
shown in \cref{eq:hyper-cross}; the successor algebra should be discharged by
the checked semiring/power normalizer in \code{IPRADerivations}, not by a large
handwritten associativity proof.

The supplied \code{lean/FiniteCore.lean} gives host definitions and the easy
telescope.  The Python checker tests 151 exact instances, but those tests are
not substitutes for \cref{eq:lean-sf-target,eq:lean-hyp-target,eq:lean-block-target}
as object derivations.

\subsection{Stage 2: rational and complex balls}

Reuse the repository's integer, rational, pairing and bounded-division
infrastructure.  A normalized fraction type is optional; an object theorem can
use raw signed numerators and positive denominators because every comparison is
cross multiplied.

Define a complex rational center as four naturals and a sign convention.  A
ball is
\[
 B(c,p)=\{z:\|z-c\|_1\le2^{-p}\}.
\]
The object predicate \code{BallEq(code,center,p)} says that the evaluator with
the fixed code returns a center whose mathematical error theorem is at most
\(2^{-p}\).  It is not necessary to encode the set braces; use a graph formula
for the evaluator and a rational inequality for the error.

Required generic lemmas are:
\begin{enumerate}
\item addition, negation, conjugation and multiplication of balls;
\item Horner evaluation with backward precision demands;
\item finite sum and product telescopes;
\item reciprocal under an apartness certificate;
\item exponential and positive logarithm Taylor evaluators;
\item positive square root by dyadic bisection;
\item path logarithm from local atanh increments.
\end{enumerate}
The Gamma packet already sketches most of these and defines host-level rational
complex arithmetic.  Each pending contract must be replaced by a bounded
checker plus a \code{Prf PRAAx} soundness theorem.

\subsection{Stage 3: named constants}

Implement fixed primitive-recursive names for:
\[
 \pi,\quad \log(2\pi),\quad \gamma,\quad
 \zeta(m)\ (m\ge2\text{ fixed}),\quad
 A,\quad \mathcal K.
\]
Recommended definitions are:
\begin{itemize}
\item Machin-type arctangent series for \(\pi\), with a geometric tail;
\item harmonic-minus-log for \(\gamma\);
\item the positive-term zeta series with integral tail;
\item the hyperfactorial Euler--Maclaurin limit \cref{eq:A-limit} for \(A\);
\item the alternating odd-square series \cref{eq:Catalan} for \(\mathcal K\).
\end{itemize}
Every name returns a rational center at precision \(p\).  Regularity is proved
by the same tail estimate that proves correctness, so no choice principle is
involved.

\subsection{Stage 4: canonical factor and product}

Define a finite rational-ball approximation to
\[
 q_k(z)=k\Log(1+z/k)-z+z^2/(2k)
\]
when \(k\ge2R\), and a direct polynomial/exponential approximation to
\(E_k(z)\) for the finite head.  The key object theorem is the cleared form of
\cref{eq:qtail}:
\begin{lstlisting}[language={}]
CanonicalTail(R,N,p) :
  2*R <= N ->
  2*R^3*2^(p+2) <= 3*N ->
  LogTailBall(R,N,p)
\end{lstlisting}
Here \code{LogTailBall} is a quantifier-free formula about the center and radius
returned by the tail evaluator.  Prove the geometric-series bound by a bounded
sum, and the \(\sum k^{-2}\) bound by the telescope
\[
 k^{-2}\le(k-1)^{-1}-k^{-1}.
\]
All divisions are cleared before induction.

Define the Barnes evaluator by
\begin{enumerate}
\item selecting \(N(R,p)\) as in \cref{eq:product-cutoff};
\item evaluating the prefactor and head factors with allocated guard bits;
\item applying the tail exponential ball;
\item multiplying the balls by a balanced tree or sequential telescope.
\end{enumerate}
A balanced tree improves numerical radii but is not logically necessary.
Soundness is an invariant over the finite factor list.

For holomorphy, implement the derivative bounds
\cref{eq:qprime-bound,eq:qsecond-bound,eq:qhigher-bound}.  It is enough initially
to formalize derivatives through order three, because all source-page theorems
can be derived from those plus fixed-order asymptotic contracts.  A later
generic order parameter can use simultaneous primitive recursion.

\subsection{Stage 5: functional equation and zeros}

The finite cancellation formula \cref{eq:finite-ratio} is a rational/power
identity after branches are exponentiated.  Reuse four Gamma contracts:

\begin{enumerate}
\item the Gauss finite product with a uniform remainder;
\item harmonic-minus-log for \(\gamma\);
\item Stirling for \((N-1)!\);
\item \((1+z/N)^N\to\e^z\) on a compact ball.
\end{enumerate}
A single telescope proves the ratio remnant \cref{eq:ratio-remnant} lies in a
ball about one.  The primary object target is
\begin{lstlisting}[language={}]
Safe(zCode,p) ->
BallEq( G(z+1), Gamma(z) * G(z), p )
\end{lstlisting}
with a local factorized version at poles.  Do not state a quotient theorem as
the primitive result.

For a fixed integer \(m\), factor the \(m\)-th canonical head polynomial as
\((z+m)^mU(z)\).  A bounded rational-box search returns a lower bound for the
other head factors, and the tail exponential is automatically apart from zero.
This produces the zero-order theorem.  Exact nonpositive-integer values are
then numeral consequences.

\subsection{Stage 6: series, integrals, and reflection}

The Taylor module uses the finite rectangle
\[
 \sum_{k\le K}\sum_{m\le M}
 \frac{(-1)^{m+1}z^m}{mk^{m-1}}
 =\sum_{m\le M}\frac{(-1)^{m+1}z^m}{m}
 \sum_{k\le K}k^{1-m}.
\]
This is a bounded-sum interchange theorem.  The two explicit tails prove
\cref{eq:taylor-tail}.  The object checker for a rational
\(\rho=a/b\) verifies \cref{eq:taylor-tail-cleared}.

For integration, reuse a generic rectangle/segment quadrature module with a
modulus field.  The log-Gamma identity is easiest to formalize from finite
products and exact antiderivatives rather than from a general theorem saying
two derivatives agree.  The improper integral needs three regions and the
regularizer described in \cref{thm:improper}.

For reflection, prove a generic finite lemma:
\begin{quote}
If a named function has a derivative ball contained in \(B(0,p+c)\) on every
cell of a certified path and its value at the initial point is zero, then its
endpoint value lies in \(B(0,p)\).
\end{quote}
Instantiate it with the cross-multiplied derivative of the Kinkelin difference.
The Clausen function is a separate sine-series name with a \(1/N\) tail.

\subsection{Stage 7: multiplication, asymptotic, and special values}

The finite product shift \cref{eq:Pshift} is a double-list cancellation.
Combine it with the already formalized Gauss multiplication theorem.  Define
the ratio \(Q_n\) only on apartness domains, but retain a cross-multiplied
periodicity theorem globally.  A right shift by a sufficiently large integer
moves every input into the asymptotic sector.  The fixed-order asymptotic
checker determines \(K(n)\), proving multiplication.

For the asymptotic module, formalize a fixed external \(N\) first.  Define the
Bernoulli polynomial \(B_{2N+2}(x)\) by a finite rational coefficient list.
Implement the periodic-Bernoulli remainder integral and prove the geometric
sector bound using a rational lower bound on \(\cos(\theta/2)\).  This suffices
for \cref{eq:Nemes-bound}.  A second theorem combines it with Gamma Stirling to
obtain \cref{eq:combined-bound}.

Special values are then short compositions.  Use a small exponent-vector
normalizer for monomials in the named positive constants
\(2,\pi,A,\Gamma(1/4),\e^{1/32},\e^{\mathcal K/(4\pi)}\).  This avoids fragile
manual rewriting of rational exponents.  The integer coefficient \(c_k\) is an
exact \code{PR} function and its half-integer theorem is qf induction.

\subsection{Stage 8: coded characterizations}

A candidate for the Vign\'eras theorem is a finite tuple of fixed programs:
\[
 (f,\log f,g',g'',g''',\text{moduli},\text{positivity bounds},
   \text{recurrence certificates}).
\]
The theorem is parameterized externally by those fixed codes.  Its premises
are object formulas saying their checkers succeed on the current finite query.
The proof in \cref{thm:vigneras} uses only right shifts, the explicit decay
bound \cref{eq:Lthird-decay}, one-period integrals, and continuity moduli.

The asymptotic characterization replaces derivative data by a relative-error
producer.  For each fixed rational \(x>0\) and precision \(p\), choose a common
\(n\) from the candidate and canonical moduli.  Periodicity then turns an
asymptotic ball into an endpoint equality ball.

\subsection{Audit and trust checks}

Each completed Lean module should finish with:
\begin{enumerate}
\item \code{\#print axioms} for host helper theorems;
\item a repository-specific linter confirming that the final theorem is a
\code{Prf PRAAx} derivation;
\item no \code{sorry}, \code{admit}, \code{axiom}, \code{Classical.choice}, or
\code{native\_decide} in the trusted path;
\item executable small-n regression examples kept logically separate from the
proof;
\item a provenance comment naming the mathematical estimate and its exact
finite version.
\end{enumerate}

The trusted endpoint is the small proof checker for the object calculus plus
Lean's kernel checking the derivation datatype.  Numerical libraries and the
Python scripts in this archive are validation aids only.

\subsection{Milestones}

A realistic sequence of independently valuable milestones is:

\begin{enumerate}
\item exact integer values, hyperfactorial relation and finite Gamma product in
\code{Prf PRAAx};
\item canonical tail and an evaluator for \(|z|\le1/2\);
\item normalization, Taylor series and the functional equation on a positive
rational domain;
\item extension to compact complex disks and zero multiplicities;
\item reflection and special half/quarter values;
\item fixed-order positive-real asymptotics;
\item fixed-sector asymptotics and multiplication;
\item the coded uniqueness theorems.
\end{enumerate}

At each stage the theorem statements remain meaningful even before the later
modules exist.  This prevents a monolithic development from hiding which
claims have actually reached object iPRA.
